\section{Step 44 Accumulate the first signed dot products}\label{h263-transform-s04}
For each input row and each basis row, multiply eight signed halfwords and add
the products. In scalar mode use full-width products and sums. In MSA mode,
\texttt{dotp\_s.w} combines pairs into four signed words; widen and reduce those
four words before storing the row result.

The input and basis signs both matter. An unsigned dot product can produce a
large positive result from a negative residual or basis constant. The final
value fits a signed word by the width proof, but the computation still needs
all eight products and their signs.

\paragraph{Observe} A negative constant block has a negative DC coefficient.
The MSA first-pass sum agrees with the scalar sum before any descaling.
\Needspace{8\baselineskip}
\paragraph{Check} Why does the signed form of the dot product matter?
\begin{itemize}
  \item[A.] It controls the memory byte order.
  \item[B.] Both samples and basis constants can be negative.
  \item[C.] It automatically quantizes the coefficients.
\end{itemize}
